% titel: Blatt 6, A1: Komplexe Zahlen
% Teil II: der durchgerechnete Weg.

\subsection*{Gegeben und gesucht}
Polarkoordinaten von $z=1+i\sqrt3$ angeben und $(1+i\sqrt3)^3$ zeigen (a);
Lösungsmenge von $|z-2i|<2$ (b); dritter Eckpunkt des gleichseitigen Dreiecks
aus $z_1=1+i\sqrt3$, $z_2=1-i\sqrt3$ (c).

\subsection*{a) Polarkoordinaten und dritte Potenz}
Für $z=1+i\sqrt3$:
\begin{equation*}
  r=\sqrt{1^2+\sqrt3^2}=2,\qquad
  \varphi=\arctan\frac{\sqrt3}{1}=\frac\pi3.
\end{equation*}
Also $z=2(\cos\frac\pi3+i\sin\frac\pi3)=2\mathrm e^{i\pi/3}$. Dann ist
\begin{equation*}
  z^3=2^3\,\mathrm e^{i\cdot3\frac\pi3}=8\,\mathrm e^{i\pi}=-8,
\end{equation*}
wie gefordert.

\subsection*{b) Lösungsmenge von $|z-2i|<2$}
$|z-2i|$ ist der Abstand von $z$ zum Punkt $2i=(0,2)$ in der komplexen Ebene.
$|z-2i|<2$ ist daher die offene Kreisscheibe mit Mittelpunkt $2i$ und Radius
$2$ (ohne Rand).

\subsection*{c) Dritter Eckpunkt des gleichseitigen Dreiecks}
$z_1,z_2$ liegen symmetrisch zur reellen Achse:
\begin{equation*}
  z_1=r\mathrm e^{i\pi/3},\quad z_2=r\mathrm e^{-i\pi/3},\qquad r=2.
\end{equation*}
Die Strecke $z_1z_2$ hat Länge $2\sqrt3$ (Differenz $2\sqrt3\,i$). Der dritte
Punkt $z_3$ liegt auf der Mittelsenkrechten, also auf der reellen Achse. Bei
einem gleichseitigen Dreieck steht die Höhe senkrecht auf der Seite; mit der
Seitenlänge $s=2\sqrt3$ und Höhe $h=\frac{\sqrt3}{2}s$ gilt
\begin{equation*}
  h=\frac{\sqrt3}{2}\cdot2\sqrt3=3.
\end{equation*}
Der Mittelpunkt von $z_1,z_2$ ist $1$. Also liefert $z_3=1\pm3$, d.h.
$z_3=4$ oder $z_3=-2$.

\begin{ergebnis}
\begin{equation*}
  \text{a)}\ r=2,\ \varphi=\tfrac\pi3,\qquad (1+i\sqrt3)^3=-8; \qquad
  \text{b)}\ |z-2i|<2\ \text{offene Scheibe bei }2i; \qquad
  \text{c)}\ z_3=4\ \text{oder}\ z_3=-2 .
\end{equation*}
\end{ergebnis}

\begin{kontrolle}
\textbf{Probe c):} Abstand $|4-(1+i\sqrt3)|=|3-i\sqrt3|=\sqrt{9+3}=2\sqrt3$,
ebenso $|4-(1-i\sqrt3)|=\sqrt{9+3}=2\sqrt3$; alle drei Seiten sind gleich
lang. Analog für $z_3=-2$: $|-2-(1+i\sqrt3)|=|-3-i\sqrt3|=\sqrt{9+3}
=2\sqrt3$.

\textbf{Probe a):} $(1+i\sqrt3)^3$ ausmultiplizieren: $(1+i\sqrt3)^2
=-2+2i\sqrt3$; mal $(1+i\sqrt3)$ ergibt $-8$ (Real- und Imaginärteil
verschwinden). -- bestätigt.
\end{kontrolle}

\begin{merksatz}
Jede komplexe Zahl als $r\mathrm e^{i\varphi}$ mit $r=|z|$ und $e^{i\varphi}$
im passenden Quadranten schreiben; Potenzen werden zu
$(r\mathrm e^{i\varphi})^n=r^n\mathrm e^{in\varphi}$. Abstände $|z-z_0|$
zeichnen Kreise; der dritte Eckpunkt eines gleichseitigen Dreiecks liegt auf
der Mittelsenkrechten des Seitenpaares.
\end{merksatz}
